\begin{frame}{Why Vector Database?}
	\begin{columns}
		\column{0.52\textwidth}
		\begin{block}{Ba giới hạn của LLM độc lập}
			\begin{itemize}
				\item \textbf{Knowledge Cutoff}: không biết dữ liệu sau thời điểm training.
				\item \textbf{Hallucination}: sinh câu trả lời có vẻ đúng nhưng thiếu evidence.
				\item \textbf{No Private Data}: không tự truy cập tài liệu nội bộ, giáo trình, log, knowledge base.
			\end{itemize}
		\end{block}
		\column{0.48\textwidth}
		\begin{alertblock}{RAG cần semantic memory}
			Vector Database lưu \textbf{embedding} của tài liệu, lập chỉ mục ANN và trả về \textbf{Top-K chunks} làm context cho LLM.
		\end{alertblock}
		\begin{exampleblock}{Khác SQL/NoSQL}
			SQL/NoSQL mạnh ở exact match; Vector Database mạnh ở \textbf{Similarity Search} trên high-dimensional vectors.
		\end{exampleblock}
	\end{columns}
	\note{Nhấn mạnh Vector Database không thay thế database truyền thống, mà giải bài toán khác: tìm theo ngữ nghĩa. Với RAG, chất lượng context phụ thuộc trực tiếp vào retrieval layer; nếu retrieval sai, LLM sẽ sinh câu trả lời sai dù model rất mạnh.}
\end{frame}

\begin{frame}{RAG Pipeline Architecture}
	\begin{columns}
		\column{0.48\textwidth}
		\begin{block}{Indexing / Ingestion}
			\begin{enumerate}
				\item Documents $\rightarrow$ Chunking.
				\item Chunks $\rightarrow$ Embedding Model.
				\item Embeddings + metadata $\rightarrow$ Vector Store.
			\end{enumerate}
		\end{block}
		\begin{block}{Retrieval / Generation}
			\begin{enumerate}
				\item User Query $\rightarrow$ Query Vector.
				\item ANN Search $\rightarrow$ Top-K relevant chunks.
				\item Context Augmentation $\rightarrow$ LLM Answer.
			\end{enumerate}
		\end{block}
		\column{0.52\textwidth}
		\begin{figure}
			\centering
			\fbox{\begin{minipage}{0.92\linewidth}
				\centering
				User Query $\rightarrow$ Embedding\\
				$\downarrow$\\
				Vector Database\\
				$\downarrow$\\
				Top-K Context Chunks\\
				$\downarrow$\\
				LLM + Prompt $\rightarrow$ Answer
			\end{minipage}}
			\caption{Workflow RAG từ query đến generated answer.}
		\end{figure}
	\end{columns}
	\note{Giải thích hai pha của RAG. Ingestion biến dữ liệu thành embedding và metadata. Retrieval biến câu hỏi thành query vector, tìm nearest neighbors rồi đưa context vào prompt. Benchmark Vector Database phải đo cả tốc độ search lẫn độ đúng của chunks được lấy ra.}
\end{frame}

\begin{frame}{The Big Three: Qdrant, Weaviate, Milvus}
	\begin{table}
		\centering
		\caption{Ba triết lý thiết kế khác nhau trong cùng bài toán Vector Search.}
		\resizebox{\textwidth}{!}{%
		\begin{tabular}{lccc}
			\hline
			\textbf{Tiêu chí} & \textbf{Qdrant} & \textbf{Weaviate} & \textbf{Milvus} \\
			\hline
			Năm ra đời & 2021 & 2019 & 2019 \\
			Core language & Rust & Go & Go + C++ \\
			Tổ chức & Qdrant Solutions & Weaviate B.V. & Zilliz + LF AI \& Data \\
			Triết lý & Performance-first & AI-native all-in-one & Distributed-first \\
			Điểm mạnh & Filtered Search & Native Hybrid Search & Billion-scale storage \\
			\hline
		\end{tabular}}
	\end{table}
	\begin{block}{Luận điểm chính}
		Qdrant tối ưu filtered vector retrieval, Weaviate tối ưu RAG workflow với Hybrid Search, Milvus tối ưu scale và High Availability.
	\end{block}
	\note{Dùng bảng này để định vị ba hệ thống trước khi đi sâu. Điểm cần tránh là xem chúng như ba sản phẩm giống hệt nhau. Chúng cùng hỗ trợ vector indexing, nhưng khác nhau ở storage model, query planner, hybrid capability và vận hành production.}
\end{frame}

\begin{frame}{Popularity \& GitHub Community}
	\begin{table}
		\centering
		\caption{Community snapshot theo nguồn tổng hợp tháng 05/2026.}
		\resizebox{\textwidth}{!}{%
		\begin{tabular}{lccc}
			\hline
			\textbf{Tiêu chí} & \textbf{Qdrant} & \textbf{Weaviate} & \textbf{Milvus} \\
			\hline
			GitHub repository & qdrant/qdrant & weaviate/weaviate & milvus-io/milvus \\
			GitHub stars & $\sim$31.4K+ & $\sim$16.2K+ & $\sim$44.4K+ \\
			Community trend & tăng nhanh từ wave RAG & ổn định, DX tốt & enterprise lớn, lâu đời \\
			Enterprise signal & startup AI agents & prototype/RAG app & Alibaba, eBay, Nvidia ecosystem \\
			\hline
		\end{tabular}}
	\end{table}
	\begin{alertblock}{Diễn giải đúng}
		GitHub stars là proxy cho cộng đồng và tài liệu, không phải thước đo tuyệt đối cho latency, accuracy hay vận hành.
	\end{alertblock}
	\note{Milvus dẫn đầu về stars nhờ lịch sử lâu hơn và backing từ Linux Foundation ecosystem. Qdrant tăng nhanh do nhu cầu RAG/agentic AI và appeal của Rust. Weaviate có cộng đồng nhỏ hơn nhưng nổi bật về developer experience và Hybrid Search.}
\end{frame}

\begin{frame}{Open-Source \& License Landscape}
	\begin{columns}
		\column{0.58\textwidth}
		\begin{table}
			\centering
			\caption{License và mô hình vận hành.}
			\scriptsize
			\begin{tabular}{lccc}
				\hline
				& \textbf{Qdrant} & \textbf{Weaviate} & \textbf{Milvus} \\
				\hline
				License & Apache 2.0 & BSD-3 & Apache 2.0 \\
				Self-hosted & Docker/K8s & Docker/K8s/Helm & Docker/K8s/Operator \\
				Managed Cloud & Qdrant Cloud & Weaviate Cloud & Zilliz Cloud \\
				Governance & Company-led & Company-led & LF AI \& Data \\
				\hline
			\end{tabular}
		\end{table}
		\column{0.42\textwidth}
		\begin{block}{Permissive OSS}
			Cả ba license đều phù hợp commercial usage và self-hosted deployment.
		\end{block}
		\begin{alertblock}{Điểm khác biệt}
			Milvus có governance độc lập hơn nhờ LF AI \& Data; Qdrant và Weaviate do công ty sáng lập dẫn dắt.
		\end{alertblock}
	\end{columns}
	\note{Không nên chỉ nhìn license. Trong production, chi phí thực tế đến từ ops complexity, dependency, backup, monitoring và khả năng mở rộng. License cho phép dùng rộng rãi, nhưng mỗi database có profile vận hành khác nhau.}
\end{frame}
